\documentclass[10pt]{amsart}
\usepackage[a4paper,margin=23mm]{geometry}
\usepackage{amsmath,amssymb,amsthm,mathtools,hyperref,graphicx,xspace,letltxmacro,xpatch}
\usepackage[english]{babel}
\usepackage{xfrac}
\providecommand{\chapterautorefname}{Chapter}

\RequirePackage{cancel}

\usepackage[scr=rsfs,cal=euler]{mathalfa}
\multlinegap0pt
\hyphenation{action actions alge-braic approach behav-ior behav-iour coho-mo-lo-gy coor-di-nate coor-di-nates define defined defines exam-ple exam-ples expect-ed explain geom-e-try inter-pre-ta-tion intro-duc-ing local-ly notion obtain obtained replac-ing sub-field togeth-er}

\newcommand{\Changel}{\mathcode`l="7160}
\newcommand{\Changelback}{\mathcode`l="716C}
\newcommand{\NoChangel}[1]{%
\expandafter\let\csname old\string#1\endcsname=#1
\let#1=\relax
\newcommand{#1}{\mathop{\mathcode`l="716C\csname old\string#1\endcsname\mathcode`l="7160}\displaylimits}%
}
\newcommand{\NoChangelnolimits}[1]{%
\expandafter\let\csname old\string#1\endcsname=#1
\let#1=\relax
\newcommand{#1}{\mathop{\mathcode`l="716C\csname old\string#1\endcsname\mathcode`l="7160}\nolimits}%
}

\NoChangelnolimits{\log}
\NoChangelnolimits{\ln}
\NoChangel{\lim}
\NoChangel{\limsup}
\NoChangel{\liminf}
\NoChangel{\varlimsup}
\NoChangel{\varliminf}
\NoChangel{\varinjlim}
\NoChangel{\varprojlim}

\let\oldbullet\bullet
\renewcommand{\bullet}{{\raisebox{1pt}{$\scriptscriptstyle\oldbullet$}}}
\let\tst\textstyle
\let\oldbigsqcup\bigsqcup
\renewcommand{\bigsqcup}{\mathop{\textstyle\oldbigsqcup}\displaylimits}
\newcommand\mto{\mathchoice{\longmapsto}{\mapsto}{\mapsto}{\mapsto}}
\newcommand\hto{\mathrel{\lhook\joinrel\to}}
\newcommand\from{\mathchoice{\longleftarrow}{\leftarrow}{\leftarrow}{\leftarrow}}
\newcommand{\To}[1]{\mathchoice{\xrightarrow{\textstyle\kern4pt#1\kern3pt}}{\stackrel{#1}{\longrightarrow}}{}{}}
\renewcommand{\labelitemi}{$\bullet$}
\newcommand{\Bset}[2]{\Bigl\{#1\Bigm|#2\Bigr\}}
\newcommand{\bbset}[2]{\biggl\{#1\Bigm|#2 \biggr\}}
\newcommand{\rbirac}[1]{\bigl(#1\bigr)}
\newcommand{\rBrac}[1]{\Bigl(#1\Bigr)}
\newcommand{\rbbrac}[1]{\biggl(#1\biggr)}
\newcommand{\cBrac}[1]{\Bigl\{#1\Bigr\}}
\newcommand{\cbbrac}[1]{\biggl\{#1\biggl\}}
\newcommand{\cbirac}[1]{\bigl\{#1\bigr\}}
\newcommand{\bcoH}[2][\bullet]{\mathsf{H}^{#1}\rbirac{#2}}
\newcommand{\bEnd}[2][]{\mathsf{End}_{#1}\rbirac{#2}}
\newcommand{\bHH}[2][\bullet]{\mathbb{H}^{#1}\rbirac{#2}}
\newcommand{\dslog}[1]{\sfrac{\dd#1}{#1}}
\newcommand{\spfrac}[2]{\sfrac{#1}{(#2)}}


\usepackage{tikz-cd}
\usepackage{tikz}
\usetikzlibrary{decorations.markings,shapes.misc,plotmarks}
\usepackage{caption}
\usepackage{subcaption}
\captionsetup[subfigure]{margin=0pt, parskip=0pt, hangindent=0pt, indention=0pt, labelformat=simple, labelfont=rm}
\renewcommand\thesubfigure{(\alph{subfigure})}

\usepackage{aliascnt}
\usepackage{xargs,ifthen}

\newcommand{\mapdef}[5]{\arraycolsep2pt
\begin{array}{crcl}
#1 \colon & #2 &\dpl\to & #3 \\[3pt]
& #4 &\dpl\mto& #5
\end{array}
}

\usepackage{bigints}
\AtBeginDocument{%
\renewcommand{\chapterautorefname}{Chapter}
\renewcommand{\sectionautorefname}{Section}
\renewcommand{\subsectionautorefname}{Section}
\renewcommand{\subsubsectionautorefname}{Section}
\renewcommand{\figureautorefname}{Figure}
\renewcommand{\itemautorefname}{Item}
\def\itemautorefname{claim}
}

\newcommandx{\thdef}[2]{
	\newaliascnt{#1}{theorem}  
	\newtheorem{#1}[#1]{#2}
	\aliascntresetthe{#1}  
	\expandafter\newcommand\expandafter{\csname #1autorefname\endcsname}{#2}
}


\newtheorem{theorem}{Theorem}[section]

\thdef{lemma}{Lemma}
\thdef{corollary}{Corollary}
\thdef{conjecture}{Conjecture}
\thdef{proposition}{Proposition}
\thdef{theoremdefn}{Theorem/Definition}

\theoremstyle{definition}
\thdef{definition}{Definition}

\theoremstyle{remark}
\thdef{remarkx}{Remark}
\thdef{examplex}{Example}

\newenvironment{example}
  {\pushQED{\qed}\renewcommand{\qedsymbol}{$\lozenge$}\examplex}
  {\popQED\endexamplex}

\newenvironment{remark}
  {\pushQED{\qed}\renewcommand{\qedsymbol}{$\lozenge$}\remarkx}
  {\popQED\endremarkx}
  
\let\defn\emph

\newcommand{\CC}{\mathbb{C}}
\newcommand{\KK}{\mathbb{K}}
\newcommand{\RR}{\mathbb{R}}
\newcommand{\QQ}{\mathbb{Q}}
\newcommand{\NN}{\mathbb{N}}
\newcommand{\ZZ}{\mathbb{Z}}
\newcommand{\AF}{\mathbb{A}}
\newcommand{\PP}{\mathbb{P}}

\newcommand{\model}[2]{\halfspc{#1}\times[0)^{#2}}
\newcommand{\halfspc}[1]{[0,\infty)^{#1}}

\newcommand{\logcyl}{\AF^1 \log\, \{0\}}
\newcommand{\logcirc}{\{0\} \log\, \{0\}}

\newcommand{\id}[1]{\mathrm{id}_{#1}}

\renewcommand{\le}{\leqslant}
\renewcommand{\ge}{\geqslant}

\newcommand{\iu}{\mathrm{i}}
\newcommand{\ipi}{\pi\iu }
\newcommand{\tipi}{2\ipi}
\newcommand{\ulog}{\mathop{\widetilde{\log}}\nolimits}
\newcommand{\uf}{\widetilde{f}}
\newcommand{\ug}{\widetilde{g}}
\newcommand{\tp}{\tilde p}
\newcommand{\ts}{\tilde s}
\newcommand{\uomega}{\underline{\omega}}

\newcommand{\defas}{\mathrel{\mathop:}=}

\newcommand{\zb}{\overline{z}}
\newcommand{\tf}{\widetilde{f}}

\let\epsilon\varepsilon

\newcommand{\set}[2]{\left\{#1\,\middle|\,#2\right\}}
\newcommand{\rbrac}[1]{\left(#1\right)}
\newcommand{\abrac}[1]{\left\langle#1\right\rangle}
\newcommand{\sbrac}[1]{\left[#1\right]}
\newcommand{\cbrac}[1]{\left\{#1\right\}}

\newcommand{\cO}[1]{\mathcal{O}_{#1}}
\newcommand{\cOx}[1]{\cO{#1}^\times}

\newcommand{\cM}{\mathscr{M}}
\newcommand{\cMgp}{\mathscr{M}^{\mathrm{gp}}}
\newcommand{\cMbas}{\mathscr{M}^{\mathrm{bas}}}
\newcommand{\cMphan}{\mathscr{M}^{\mathrm{phan}}}

\newcommand{\Cinf}[1]{\mathscr{C}^{\infty}_{#1}}
\newcommand{\Cinfpos}[1]{\mathscr{C}^{\infty, > 0}_{#1}}
\newcommand{\Cinfge}[1]{\mathscr{C}^{\infty, \ge 0}_{#1}}
\newcommand{\Cinflog}[1]{\mathscr{C}^{\infty,\log}_{#1}}

\newcommand{\dd}{\mathrm{d}}
\newcommand{\dlog}[1]{\frac{\dd#1}{#1}}
\newcommand{\cA}[2][\bullet]{\mathscr{A}^{#1}_{#2}}\newcommand{\A}[2][\bullet] {\mathscr{A}^{#1}\rbrac{#2}}
\newcommand{\Ac}[2][\bullet]{\mathscr{A}_c^{#1}\rbrac{#2}}
\newcommand{\cAlog}[2][\bullet]{\mathscr{A}^{#1,\log}_{#2}}
\newcommand{\Alog}[2][\bullet]{\mathscr{A}^{#1,\log}(#2)}
\newcommand{\Alogc}[2][\bullet]{\mathscr{A}^{#1,\log}_c(#2)}

\newcommand{\VFlog}[1]{\mathcal{T}^{\log}_{#1}}
 \newcommand{\coVFlog}[1]{\mathcal{T}^{\vee,\log}_{#1}}
\newcommand{\cvf}[1]{\partial_{#1}}
\newcommand{\logcvf}[1]{#1\partial_{#1}}

\newcommand{\reg}[1][]{\mathrm{reg}_{#1}}
\DeclareMathOperator{\reglim}{reglim}
\DeclareMathOperator*{\res}{Res}

\newcommand{\Sig}{\Sigma}
\newcommand{\uSig}{\underline{\Sigma}}
\newcommand{\uSigo}{\underline{\Sigma}^\circ}
\newcommand{\uphi}{\underline{\phi}}
\newcommand{\uPsi}{\underline{\Psi}}
\newcommand{\bas}[1]{{#1}^{\mathrm{bas}}}
\newcommand{\phant}{\mathrm{phan}}
\newcommand{\phan}[1]{{#1}^{\phant}}
\newcommand{\Sigbas}{\bas{\Sig}}
\newcommand{\ubdSig}[1][]{\underline{\partial^{#1}\Sigma}}
\newcommand{\bdSig}[1][]{\partial^{#1}\Sigma}
\newcommand{\bdtopSig}[1][]{\partial^{#1}_{\mathrm{top}}\Sigma}
\newcommand{\bdSigbas}[1][]{\partial^{#1}\Sigbas}

\newcommand{\intSig}{\Sigma^\circ}
\newcommand{\uintSig}{\underline{\Sigma}^\circ}

\newcommand{\intPsi}{\Psi^\circ}
\newcommand{\Psibas}{\bas{\Psi}}

\newcommand{\tb}[2][]{T_{#1}#2}
\newcommand{\tbphan}[2][]{T_{#1}^{\phant}#2}
\newcommand{\tbpos}[2][]{T^{>0}_{#1}#2}
\newcommand{\tbposphan}[2][]{T^{\phant,>0}_{#1}#2}
\newcommand{\tbge}[2][]{T^{\ge0}_{#1}#2}
\newcommand{\tbgephan}[2][]{T^{\phant,\ge0}_{#1}#2}
\newcommand{\ctb}[2][]{T^\vee_{#1}#2}

\newcommand{\cbdX}{\overline{\partial X}}
\newcommand{\bdX}[1][]{\partial^{#1} X}
\newcommand{\cX}{\overline{X}}
\newcommand{\uX}{\underline{X}}
\newcommand{\uY}{\underline{Y}}
\newcommand{\cY}{\overline{Y}}
\newcommand{\bD}{\overline{D}}
\newcommand{\uD}{\underline{D}}
\newcommand{\bA}{\overline{A}}
\newcommand{\bB}{\overline{B}}
\newcommand{\KN}[1]{\mathrm{KN}(#1)}
\newcommand{\KNR}[1]{\mathrm{KN}_\RR\rbrac{#1}}
\newcommand{\Spec}[1]{\mathsf{Spec}(#1)}
\newcommand{\Dbl}[1]{\mathrm{Dbl}(#1)}

\newcommand{\sect}[1]{\mathsf{\Gamma}(#1)}
\newcommand{\coH}[2][\bullet]{\mathsf{H}^{#1}\rbrac{#2}}
\newcommand{\coHc}[2][\bullet]{\mathsf{H}^{#1}_{\mathrm{c}}\rbrac{#2}}
\newcommand{\HH}[2][\bullet]{\mathbb{H}^{#1}\rbrac{#2}}
\newcommand{\HdR}[2][\bullet]{\mathsf{H}_{\mathrm{dR}}^{#1}(#2)}
\newcommand{\HdRc}[2][\bullet]{\mathsf{H}_{\mathrm{dR},\mathrm{c}}^{#1}(#2)}
\newcommand{\HB}[2][\bullet]{\mathsf{H}_{\mathrm{B}}^{#1}(#2)}
\newcommand{\HlgyB}[2][\bullet]{\mathsf{H}^{\mathrm{B}}_{#1}(#2)}

\newcommand{\Hlgy}[2][\bullet]{\mathsf{H}_{#1}(#2)}

\newcommand{\forms}[2][\bullet]{\Omega^{#1}_{#2}}

\newcommand{\End}[2][]{\mathsf{End}_{#1}\rbrac{#2}}
\newcommand{\GL}[2]{\mathrm{GL}_{#1}(#2)}
\newcommand{\symgp}[1]{\mathfrak{S}_{#1}}
\newcommand{\sgn}[1]{\mathrm{sgn}_{#1}}
\newcommand{\unitcirc}{\mathsf{S}^{1}}

\newcommand{\simp}[3]
{
\begin{tikzcd}[ampersand replacement=\&]
     {#1} \& {#2} \ar[l,shift left = 1]\ar[l,shift right=1] \& {#3} \ar[l,shift left = 2]\ar[l,shift left = 0]\ar[l,shift left = -2] \& \ar[l,shift left = 3]\ar[l,shift left = 1]\ar[l,shift left = -1]\ar[l,shift left = -3] \cdots
    \end{tikzcd}
}

\newcommand{\augsimp}[4]
{
\begin{tikzcd}[ampersand replacement=\&]
     #1 \& #2 \ar[l] \& #3
     \ar[l,shift left = 1]\ar[l,shift right=1] \& #4 \ar[l,shift left = 2]\ar[l,shift left = 0]\ar[l,shift left = -2] \& \ar[l,shift left = 3]\ar[l,shift left = 1]\ar[l,shift left = -1]\ar[l,shift left = -3] \cdots
    \end{tikzcd}
}

\newcommand{\augsimpshort}[4]
{
\begin{tikzcd}[ampersand replacement=\&]
     #1 \& #2 \ar[l] \& #3
     \ar[l,shift left = 1]\ar[l,shift right=1] \& \ar[l,shift left = 2]\ar[l,shift left = 0]\ar[l,shift left = -2]  #4
    \end{tikzcd}
}

\newcommand{\scat}{\mathscr{I}}
\newcommand{\ascat}{\mathscr{I}_+}
\newcommand{\op}{\mathrm{op}}
\DeclareMathOperator*{\holim}{holim}
\DeclareMathOperator*{\fibre}{fiber}

\usepackage{xcolor}
\let\dpl\displaystyle
\providecommand{\skpt}{\par\smallskip\noindent}
\providecommand{\pointir}{.\ }
\providecommand{\eg}{e.g.}\providecommand{\ie}{i.e.}
\providecommand{\resp}{resp.}\providecommand{\cf}{cf.}\providecommand{\etc}{etc.}
\let\Subsection\subsection\let\Subsubsection\subsubsection
\AtBeginDocument{\gdef\shorttitle{Regularized integration on log corners}}

\providecommand{\mainmatter}{}
\providecommand{\backmatter}{}
